\section{Why this combination works}
\label{sec:why}

Each tool has a sweet spot. The trick is to stop asking any one of them to do what it is bad at.

\subsection{Overleaf}
Overleaf is a writing environment. Its strengths are:
\begin{itemize}[leftmargin=*]
  \item Live PDF preview with sub-second compile-on-save.
  \item Zero-install collaboration --- co-authors click a link and type.
  \item A complete TeX Live with most packages pre-installed.
  \item Rich-text mode for co-authors who have never seen \LaTeX{}.
  \item Track changes, comments, and a linear version history.
\end{itemize}

It is deliberately \emph{not} a programmer's editor. Multi-file search-and-replace is clunky, refactoring is manual, and you cannot ask it ``find every \verb|\cite| whose key is missing from the \texttt{.bib}.''

\subsection{Claude Code}
Claude Code runs locally in your terminal and has read/write access to every file in the project. For a \LaTeX{} project that means it can:
\begin{itemize}[leftmargin=*]
  \item Read all section files and the bibliography at once.
  \item Make coordinated edits across many files in a single pass.
  \item Follow conventions from a \verb|CLAUDE.md| that it picks up automatically each session.
  \item Run tools (\texttt{latexmk}, \texttt{chktex}, \texttt{latexdiff}) and react to their output.
  \item Produce reviewable diffs rather than opaque rewrites.
\end{itemize}

It is bad at watching a PDF compile in real time and it is overkill for fixing a typo in one place.

\subsection{GitHub}
GitHub is the glue. It gives you:
\begin{itemize}[leftmargin=*]
  \item A durable history of every edit, keyed by commit message.
  \item Branches for parallel work (a revision round, a response letter, an alternative introduction).
  \item Pull requests as a review surface for co-authors or for you reviewing Claude's edits.
  \item Continuous integration that can build the PDF on every push, so nothing broken reaches your co-authors.
\end{itemize}

\subsection{What you get by combining them}
\begin{itemize}[leftmargin=*]
  \item \textbf{Fast writing loop} in Overleaf --- stays in flow while drafting.
  \item \textbf{Bulk and structural edits} via Claude Code --- minutes instead of hours.
  \item \textbf{An audit trail} via GitHub --- every change is inspectable, revertible, and attributable.
  \item \textbf{A single source of truth} --- the repository --- that both tools write into.
\end{itemize}

The result is that each tool does the job it is best at, and you only ever touch the one that is right for the job in front of you.
